\bigskip

\begin{minipage}{0.41\linewidth}
Dans l'espace muni d'un repère orthonormé \Oijk, on considère les points 

\[\text{A}(0~;~8~;~6),\quad  \text{B}(6~;~4~;~4 )\quad \text{et C}(2~;~4~;~0).\]

\end{minipage}\hfill
\begin{minipage}{0.55\linewidth}
\psset{unit=1cm,arrowsize=2pt 3}
\begin{pspicture}(-1.8,-1.8)(5.4,3.8)
\pspolygon[fillstyle=solid,fillcolor=lightgray](1.7,-0.6)(4.8,3)(0.4,0.4)
\psaxes[linewidth=1.25pt,Dx=20,Dy=20]{->}(0,0)(0,0)(5.4,3.8)
\psline[linewidth=1.25pt]{->}(0,0)(-1.8,-1.8)
\multido{\n=0+-0.3}{7}{\psdots[dotstyle=+,dotangle=45](\n,\n)}
\multido{\n=0+0.6}{9}{\psline(\n,-0.05)(\n,0.05)}
\multido{\n=0+0.5}{8}{\psline(-0.05,\n)(0.05,\n)}
\psline[linestyle=dashed](4.8,0)(4.8,3)(0,3)
\psframe[linestyle=dashed](-1.8,-1.8)(0.4,0.4)
\psframe[linestyle=dashed](2.4,2)
\psline[linestyle=dashed](-1.8,0.4)(0,2)
\psline[linestyle=dashed](0.4,-1.8)(2.4,0)
\psline[linestyle=dashed](0.4,0.4)(2.4,2)
\psline[linestyle=dashed](-0.6,-0.6)(1.8,-0.6)
\uput[d](0,0){\small O} \uput[ul](-0.15,-0.2){\footnotesize $\vect{\imath}$}
\uput[d](0.3,0){\footnotesize $\vect{\jmath}$}
\uput[l](0,0.3){\footnotesize $\vect{k}$}
\uput[l](-1.8,-1.8){$x$}\uput[d](5.3,0){$y$}\uput[l](0,3.7){$z$}

\uput[ur](4.8,3){A} \uput[ul](0.4,0.4){B} \uput[dr](1.7,-0.6){C}
\end{pspicture}
\end{minipage}

\bigskip

\begin{enumerate}
\item 
	\begin{enumerate}
		\item Justifier que les points A, B et C ne sont pas alignés.
		\item Montrer que le vecteur $\vect{n}(1~;~2~;~-1)$ est
un vecteur normal au plan (ABC).
		\item Déterminer une équation cartésienne du plan (ABC).
	\end{enumerate}
\item Soient D et E les points de coordonnées respectives (0~;~0~;~6) et (6~;~6~;~0).
	\begin{enumerate}
		\item Déterminer une représentation paramétrique de la droite (DE).
		\item Montrer que le milieu I du segment [BC] appartient à la droite (DE).
	\end{enumerate}	
\item On considère le triangle ABC.
	\begin{enumerate}
		\item Déterminer la nature du triangle ABC.
		\item Calculer l'aire du triangle ABC en unité d'aire.
		\item Calculer $\vect{\text{AB}}  \cdot \vect{\text{AC}}$.
		\item En déduire une mesure de l'angle $\widehat{\text{BAC}}$ arrondie à $0,1$ degré.
	\end{enumerate}	
\item On considère le point H de coordonnées $\left(\dfrac53~;~\dfrac{10}{3}~;~- \dfrac53\right)$.

Montrer que H est le projeté orthogonal du point O sur le plan (ABC). 

En déduire la distance du point O au plan (ABC).
\end{enumerate}
